%! Logarithmic Function Context and Data Modeling — Unit 5, Lesson 5.7 — Exit Ticket
\documentclass[10pt]{article}
\usepackage{precalculus-article}
\usepackage{precalculus-boxes}
\newcommand{\TallMath}[1]{$\displaystyle #1\rule[-1.4em]{0pt}{3.2em}$}

\begin{document}

\pageheader{Unit 5, Lesson 5.7}{Exit Ticket}
\namedateperiod

\vspace{0.15in}

\begin{enumerate}[leftmargin=1.8em, itemsep=18pt]

\item The table below shows a data set.

\vspace{6pt}
\begin{center}
\renewcommand{\arraystretch}{1.8}
\begin{tabular}{|c|c|c|c|c|}
\hline
\rowcolor{sky}
$x$    & $2$ & $4$ & $8$ & $16$ \\
\hline
$f(x)$ & $1$ & $2$ & $3$ & $4$ \\
\hline
\end{tabular}
\end{center}

\begin{enumerate}[leftmargin=1.5em, itemsep=8pt, label=\alph*.]
  \item What is the proportion factor (how does $x$ change as $f(x)$ increases by 1)?

  \vspace{4pt}
  Proportion: \blank{2cm} \qquad Base $b = $ \blank{2cm}

  \item Write the logarithmic model in the form $f(x) = \log_b(x)$.

  \vspace{4pt}
  $f(x) = $ \blank{4cm}

  \item Use the model to predict $f(32)$.

  \vspace{4pt}
  $f(32) = $ \blank{4cm}
\end{enumerate}

\item Let $f(x) = 3\ln(x)$.

\begin{enumerate}[leftmargin=1.5em, itemsep=8pt, label=\alph*.]
  \item Evaluate $f(e^2)$. Show each step.

  \vspace{4pt}
  $f(e^2) = 3\ln(e^2) = 3 \cdot $ \blank{1.5cm} $= $ \blank{2cm}

  \item Evaluate $f(1)$.

  \vspace{4pt}
  $f(1) = 3\ln(1) = 3 \cdot $ \blank{1.5cm} $= $ \blank{2cm}

  \item Solve $f(x) = 6$ for $x$. Show your work.

  \vspace{4pt}
  $3\ln(x) = 6 \;\Rightarrow\; \ln(x) = $ \blank{2cm} $\;\Rightarrow\; x = $ \blank{3cm}
\end{enumerate}

\end{enumerate}

\end{document}
